\section{Definitions}
\label{sec:preliminaries}

All graphs produced by our reductions are finite, simple, and
unweighted. Source graphs may have weights, parallel edges, or loops
when explicitly stated. An \emph{independent set} contains no adjacent
pair of vertices.
For an unweighted graph $G$, let $\alphaI(G)$ denote the size of a maximum independent set.
A \emph{clique} has every pair of distinct vertices adjacent.
A \emph{vertex cover} is a subset of vertices $X$ such that each edge is incident to a vertex in $X$.
Then, $\vc(G)$ the minimum vertex-cover size.
Taking complements gives


For $U\subseteq V(G)$, we write $G[U]$ for the subgraph induced by $U$, $G[U]=(U,\{uv\in E(G)\colon uv\in V(G)\})$, $G-U$ for
$G[V(G)\setminus U]$, $\Delta(G)$ for maximum degree, and
$[q]=\{1,\ldots,q\}$. A spanning subgraph retains all vertices
and some of the edges. A \emph{matching} is a set of edges with
pairwise disjoint endpoints, and a graph is \emph{cubic} if every
vertex has degree three.
A \emph{connected component} of a graph is any of its maximal connected subgraphs.

In a rooted forest, two distinct vertices are \emph{comparable} if one is
an ancestor of the other; each vertex is also comparable with itself.
The subtree rooted at a vertex $v$ consists of $v$ and all its descendants.
For a rooted forest $F$ and $U\subseteq V(F)$, the \emph{restriction}
$F|_U$ is the rooted forest with vertex set $U$ and edge set $E(F|_U)$ consisting of all $uv$, $u,v\in U$, such that $u$ is an ancestor of $v$ in $F$, and no $w\in U$ is an internal vertex of the  $uv$-path in $F$.
(Hence, the roots in $F|_U$ are the vertices of $U$ with no ancestor in $U$).
Denote by $\operatorname{depth}_F(v)$ the \emph{depth} of $v$ in $F$ as the number of its ancestors plus one; hence roots have depth one.

An \emph{elimination forest} $F$ of a graph $G$ is a rooted forest on
$V(G)$ such that the endpoints of every edge of $G$ are comparable in $F$.
The height $h(F)$ is the maximum number of vertices on a
root-to-leaf path. The \emph{treedepth} $\td(G)$ of a graph $G$ is the minimum height of
an elimination forest of $G$. We set the treedepth to be $0$ for a graph with no vertices.
For disconnected graphs, treedepth is the maximum treedepth of a connected component.

The \emph{total depth} of an elimination forest $F$ is
$\totd(F)=\sum_{v\in V(F)}\operatorname{depth}_F(v)$.
Let $\totd(G)$ be the minimum total depth of an
elimination forest of $G$. This objective is also called
\emph{VPT-sum}~\cite{HogemoEtAl2021}.

A \emph{tree decomposition} of $G$ is a pair $(T,B)$, where $B=\{B_t\subseteq V(G):t\in V(T)\}$ is a collection of \emph{bags}, satisfying three conditions:
\begin{enumerate}
\item $\displaystyle\bigcup_{t\in V(T)} B_t=V(G)$.
\item For every $uv\in E(G)$, there is a $t\in V(T)$ such that
$\{u,v\}\subseteq B_t$.
\item For every vertex $v$ of $G$, the nodes $t$ with $v\in B_t$
induce a connected subtree of $T$.
\end{enumerate}
Its width is
$\max_t|B_t|-1$, and the \emph{treewidth} of $G$, denoted $\operatorname{tw}(G)$, is
the minimum width taken over all three decompositions.
